%! TEX root = ../main.tex

\chapter{Cosmic rays and extensive air showers}
\label{chap:cosmic-rays}

\CRs are particles of extraterrestrial origin that travel very close to the 
speed of light. Their relativistic kinetic energy $pc$ far surpasses their rest 
mass $mc^2$. In particular, \UHECRs are typically defined as \CRs with energies
exceeding $\SI{1}{\exa\eV} = \SI{e18}{\eV}$ 
\cite{alvesbatistaOpenQuestionsCosmicRay2019} These microscopic fragments 
originate from some of the most dynamic processes in the universe. As such, they
are of great interest to researchers. Not only do they hold the key to 
understanding particle interactions at extremely high energies and small scales,
but can also improve the understanding of the cosmos over the largest distances.

\section{Cosmic Ray Phenomenology}
\label{sec:cr-phenomenology}

\fulllift

\subsection{Historical Discovery}
\label{ssec:cr-history}

The study of cosmic rays is an active area of research and a field more than one
hundred years in the making. The turn of the twentieth century saw many 
physicists testing the predictions of the newly developed theory of 
electromagnetism. Of special importance for this story are the advancements of 
Theodor Wulf, who, in 1909, developed an electrometer that is able to measure 
the ionization rate of the atmosphere. His invention was used to show the 
difference in air ionization rates on the ground, at the top of the Eiffel Tower
\cite{wulfBeobachtungenUberStrahlung1910} and underwater
\cite{paciniRadiazionePenetranteAlla1912}. In 1912 Victor Hess continued this
study by examining spontaneous ionization during ballon flights above Germany 
and Austria \cite{hessElektrischeLeitfaehigkeitAtmosphaere1926}. With the 
ionization rate increasing the higher the balloon flew Hess concluded that 
energetic particles are present in the upper atmospher. His intuition proved 
correct and would eventually be honored with a Nobel prize in 1936 
\cite{nobelprizeoutreach2025NobelPrizePhysicsa}.

Given the context of research at the time, these energetic particles were
thought to be \EM radiation. Their extraterrestrial origin was demonstrated by 
Hess himself, as well as Werner Kolhörster by 1914 
\cite{kolhorsterMeasurementsPenetratingRadiation1913}. Consequentially, the 
term \textit{cosmic rays} was coined for the phenomenon. In 1929 Walter Bothe 
and Kolhörster attempted to investigate the interactions of cosmic rays with 
high-$Z$ absorbers \cite{botheNaturHohenstrahlung1929}. They were surprised to
find out that the particles they investigated had a penetrating power 
uncharacteristic for $\upgamma$-radiation. They concluded that \CRs are made up 
of, at least partially, corpuscular radiation, rather than photons.

Since then, the study of \CRs has lead to the discovery of the positron 
(Anderson 1933, Ref.\ \cite{andersonPositiveElectron1933}), the muon (Anderson,
Neddermeyer 1936, Ref.\ \cite{andersonElementaryParticlesMatter1949}), and a 
plethora of other particles. Today we know that \CRs are very diverse in nature.
Their composition ranges from elementary particles like electrons, over exotic
mesons, to heavy nuclei. Photons too have been detected in \CR experiments,
proving that the original name is not a complete misnomer.

An important milestone was further achieved by Pierre Auger in 1938. He 
conducted coincidence measurements in the swiss Alps and found that the rate of
coincidences between several \CR detectors is directly proportional to their 
separation \cite{augerExtensiveCosmicRayShowers1939}. This marks the discovery 
of \EAS, and the beginning of ground based \CR research in array detectors.

\subsection{CR Observables}

The many discoveries of new fragments in the cosmic radiation has resulted in
the creation of entirely new theories that aim to describe the diverse particle 
zoo at hand. In general though, only a handful parameters, called hereafter \CR 
observables, are sufficient to characterize the kinematics of a cosmic ray from 
genesis at the source, over propagation, to detection on Earth. It follows a 
brief introduction to the most important properties.

\subsubsection{Energy}

The energy of the \CR particle is perhaps its' most defining observable. After 
a particle is ejected from a \CR acceleration mechanism (\todo{ref}) it begins
a journey mostly consisting of adiabatic processes. This means the \CR energy 
dictates how (and how far) the particle propagates through spacetime, as well as
the size of the atmospheric particle cascade it kicks off on Earth.

\todo{ref} will reveal that the energy distribution, also called the energy 
spectrum, of most \CR accelerators is a power law $E^{-\alpha}$ with a single
spectral index $\alpha$. The energy spectrum measured on Earth is not, as shown 
in \cref{fig:cr-spectrum}. It shows several different kinks, corresponding to a 
hardening (softening) of the spectrum, where the average energy of a \CR 
increases (decreases). These kinks in general correspond to transition between 
different \CR populations, that each have a respective particle flux, injection 
and cutoff energy. An example of how these transitions result in kinks is shown 
in \cref{fig:cr-spectrum-explanation}. To date, several transition have been 
observed in the \CR spectrum observed on Earth:

\begin{itemize}
	\vspace{0.7cm}
	\item \textbf{The knee} refers to a softening of the spectrum around 
	\SI{4}{\peta\eV}. The spectral index changes from $\alpha\approx2.7$
	below the knee to $\alpha\approx3.1$ above the knee 
	\cite{blumerCosmicRaysKnee2009}. It has been shown that this coincides
	with a transition in the mean mass of \CR particles 
	\cite{lhaasocollaborationMeasurementsAllParticleEnergy2024}. The knee is
	thus thought to mark the transition from light \CRs 
	($m_\mathrm{CR} \leq m_\mathrm{He}$) to heavier hadronic cosmic rays.

	\vspace{0.7cm}
	\item \textbf{The second knee} is a further softening from 
	$\alpha\approx3.0$ at an energy below roughly \SI{500}{\peta\eV} to 
	$\alpha\approx3.25$ above \cite{naganoEnergySpectrumPrimary1992, 
	lawrenceCosmicRayEnergy1991, birdEvidenceCorrelatedChanges1993, 
	abu-zayyadMeasurementCosmicRayEnergy2001}. Like the softening that 
	occurs at \SI{4}{\peta\eV}, this feature is assumed to be composition
	related. However, mass sensitive studies around the second knee to date
	have not been able to confirm this theory with sufficient certainty.

	\vspace{0.7cm}
	\item \textbf{The ankle} marks the transition from $\alpha\approx3.0$ to
	$\alpha\approx2.5$ at an energy of roughly \SI{5}{\exa\eV} 
	\cite{fenuCosmicRayEnergy2023}. The Pierre Auger Observatory has 
	reported on a dipole in the arrival directions of \CRs above 
	\SI{8}{\exa\eV} 
	\cite{thepierreaugercollaborationObservationLargescaleAnisotropy2017}, 
	coincident with the hardening of the energy spectrum. Since the 
	anisotropy points away from the galactic center the ankle is projected
	to represent the transition from a predominantly galactic \CR population
	to an extragalactic \CR population.

	\vspace{0.7cm}
	\item \textbf{The instep} is a softening of the \UHECR spectrum at 
	\SI{13}{\exa\eV} where $\alpha$ decreases from $\sim2.5$ to $\sim3.0$
	\cite{aabMeasurementCosmicrayEnergy2020}. This feature can be reproduced
	in simulations when altering the \CR mass composition from extragalactic
	light \CRs to extragalactic heavy \CRs 
	\cite{tsunesadaJointAnalysisEnergy2022}. Similar to the second knee 
	feature, experiments to date do not have a sufficiently good mass 
	sensitivity to confirm this theory.

	\vspace{0.7cm}
	\item \textbf{The suppression} marks the steep falloff at the tail-end 
	of the (measurable) \CR energy spectrum. At an energy above 
	$\sim\SI{50}{\exa\eV}$ the spectral index drastically changes to 
	$\alpha\approx5.4$ \cite{aabMeasurementCosmicrayEnergy2020}. Cosmic 
	rays exceeding this energy are incredibly sparse. It is theorized that
	no process exists in the (local) cosmos that can reliably accelerate 
	particles beyond this cutoff energy 
	\cite{tsunesadaJointAnalysisEnergy2022, greisenEndCosmicRaySpectrum1966,
	kuzminCosmicrayInteractionsPhotons1968}.
\end{itemize}

\begin{figure}[t]
  \centering
  \subfloat[]{\includegraphics[height=0.30\textwidth]{cosmic-rays/spectrum.png}
  \label{fig:cr-spectrum}
  }\hspace{0.2cm}
  \subfloat[]{\includegraphics[height=0.30\textwidth]{cosmic-rays/spectrum-explanation.png}
  \label{fig:cr-spectrum-explanation}
  }
  \caption[]{\subref{fig:cr-spectrum} The distribution of energies of \CRs 
	  arriving on earth. Several transition regions are marked and explained
	  in the main text. Data stems from 
	  \cite{re3data.orgKASCADECosmicRay2013} and 
  	  \cite{yushkovMassCompositionCosmic2021}. 
  	  \subref{fig:cr-spectrum-explanation} The observed energy spectrum 
  	  (black) is the sum of different \CR populations (grey) that cannot
 	  easily be distinguished. The cutoff energies of the individual spectra
  	  coincide with a softening of the total spectrum. The transition from 
  	  one population to the other can result in a hardening of the \CR
 	  spectrum observed on Earth.}
  \label{fig:energy-spectrum}
\end{figure}

\subsubsection{Mass}

An almost negligible fraction of the total energy is (for massive particles) 
given by the rest mass of the cosmic ray. Given the miniscule relative strength
of gravity\footnote{Compared to the \EM-, strong or weak force.} mass terms can 
be safely ignored in \CR acceleration and propagation. 


%	\item The \textbf{electromagnetic charge} of a cosmic ray defines how it
%	      interacts with external electric and magnetic fields.
%\end{enumerate}


\subsubsection{Charge}

The electromagnetic charge of a cosmic ray defines how it interacts with 
externel electric and magnetic fields. 

\subsubsection{Rigidity}



\section{Origins of Cosmic Rays}
\label{sec:cr-accelerators}

\section{Cosmic Ray Propagation}
\label{sec:cr-propagation}

\section{Extensive Air Showers}
\label{sec:extensive-air-showers}



